\section{Attribute Assignment Details}
\label{app:attribute-assignment-details}

\noindent\textbf{CelebA and ImageNet10.}
Attribute assignments for CelebA use the official CelebA annotations and follow the SafeFix setting~\citep{xu2025safefix}. ImageNet10 does not provide official attribute annotations, so we follow the SafeFix-style ImageNet10 protocol: define task-level attribute fields such as \texttt{color}, \texttt{background}, and \texttt{texture}, then assign closed-set attribute values to each image with the VLM. These assignments are used only as selection metadata for \model and for post-training slice diagnostics; benchmark labels remain the source of the reported task metrics.

\noindent\textbf{COCO and KITTI.}
For COCO pose estimation and KITTI object detection, we follow HiBug2~\citep{chen2025hibug2} and use task-specific closed-set attribute fields proposed before source selection. We take KITTI as a concrete example. GPT-5~\citep{openai2025gpt5} proposes 10 representative object-detection categories, listed in Figure~\ref{tab:kitti-attribute-corpus}. Each category has a predefined set of candidate values, a natural-language prompt template, and a corresponding query directed to Qwen3-VL~\citep{bai2025qwen3vl}. The repair fields used in the main KITTI experiments are color and weather; the remaining fields are retained for broader slice analysis and diagnostics.

\begin{figure}[htbp]
\centering
\caption{KITTI object-detection attribute corpus used for VLM assignment.}
\label{tab:kitti-attribute-corpus}
\begin{tcolorbox}[
    colback=blue!2!white,
    colframe=blue!15!black,
    title=\textbf{KITTI Attribute Corpus \& VLM Prompts},
    fonttitle=\bfseries,
    rounded corners,
    drop shadow,
    left=6pt,
    right=6pt,
    top=4pt,
    bottom=4pt,
    fontupper=\small,
    width=0.98\linewidth
]
\textit{Note: The placeholder \texttt{\#1} is replaced by a candidate value from the corresponding category.}
\vspace{0.8em}

\begin{description}
    \setlength{\itemsep}{0.6em}
    \item[\texttt{color}:] \texttt{white}, \texttt{black}, \texttt{gray}, \texttt{silver}, \texttt{red}, \texttt{blue}, \texttt{green}, \texttt{yellow}, \texttt{brown} \\
    \textcolor{gray}{\textbf{Prompt:}} \texttt{a \#1 main target object in a driving scene}

    \item[\texttt{object\_scale}:] \texttt{small}, \texttt{medium}, \texttt{large} \\
    \textcolor{gray}{\textbf{Prompt:}} \texttt{a \#1 target object in a driving scene}

    \item[\texttt{viewpoint}:] \texttt{front}, \texttt{rear}, \texttt{side}, \texttt{front-side}, \texttt{rear-side} \\
    \textcolor{gray}{\textbf{Prompt:}} \texttt{a \#1 view of the main target object in a driving scene}

    \item[\texttt{occlusion}:] \texttt{none}, \texttt{partial}, \texttt{heavy} \\
    \textcolor{gray}{\textbf{Prompt:}} \texttt{a \#1 occluded target object in a driving scene}

    \item[\texttt{truncation}:] \texttt{none}, \texttt{partial}, \texttt{heavy} \\
    \textcolor{gray}{\textbf{Prompt:}} \texttt{a \#1 truncated target object in a driving scene}

    \item[\texttt{motion\_state}:] \texttt{moving}, \texttt{stationary} \\
    \textcolor{gray}{\textbf{Prompt:}} \texttt{a \#1 target object in a driving scene}

    \item[\texttt{background\_style}:] \texttt{city street}, \texttt{residential street}, \texttt{highway}, \texttt{parking area}, \texttt{service lane}, \texttt{tree-lined road} \\
    \textcolor{gray}{\textbf{Prompt:}} \texttt{a driving scene on a \#1}

    \item[\texttt{traffic\_density}:] \texttt{sparse}, \texttt{light}, \texttt{moderate}, \texttt{heavy} \\
    \textcolor{gray}{\textbf{Prompt:}} \texttt{a driving scene with \#1 traffic}

    \item[\texttt{lighting\_condition}:] \texttt{bright}, \texttt{overcast}, \texttt{strong shadows}, \texttt{glare}, \texttt{dim} \\
    \textcolor{gray}{\textbf{Prompt:}} \texttt{a driving scene with \#1 lighting}

    \item[\texttt{weather}:] \texttt{clear}, \texttt{cloudy}, \texttt{rainy}, \texttt{foggy}, \texttt{snowy} \\
    \textcolor{gray}{\textbf{Prompt:}} \texttt{a driving scene in \#1 weather}
\end{description}
\end{tcolorbox}
\end{figure}

